\documentclass[11pt]{article}
% ===========================================================================
%  Shared preamble for the AP Calculus AB exam-preparation worksheets.
%  Input right after \documentclass. Files beginning with "_" are not compiled
%  on their own. Navy and gold, to match the course's lesson pages.
% ===========================================================================
\usepackage[letterpaper,top=2.3cm,bottom=2.2cm,left=1.7cm,right=1.7cm,headheight=28pt,footskip=30pt]{geometry}
\usepackage{amsmath,amssymb}
\usepackage[T1]{fontenc}
\usepackage[utf8]{inputenc}
\usepackage{lmodern}
\usepackage{xcolor}
\usepackage[most]{tcolorbox}
\usepackage{enumitem}
\usepackage{fancyhdr}
\usepackage{array}
\usepackage{tikz}
\usepackage{pgfplots}
\pgfplotsset{compat=1.18}

\definecolor{navy}{HTML}{161E54}
\definecolor{ink}{HTML}{27357E}
\definecolor{gold}{HTML}{B98A00}
\definecolor{goldbg}{HTML}{FFFDF3}
\definecolor{inkbg}{HTML}{F6F7FD}
\definecolor{rule}{HTML}{C9CFF0}
\definecolor{okgreen}{HTML}{1B7A44}
\definecolor{okbg}{HTML}{F1F8F3}

\pagestyle{fancy}
\fancyhf{}
\fancyhead[L]{\footnotesize Name: \underline{\hspace{4cm}}}
\fancyhead[C]{\textbf{\color{navy}Integration Academy}\\[-2pt]{\scriptsize Calculus AB \textperiodcentered{} AP\textsuperscript{\textregistered} Exam Preparation}}
\fancyfoot[L]{\tiny\color{black!55}\parbox[t]{0.78\textwidth}{AP\textsuperscript{\textregistered} is a registered trademark of the College Board, which was not involved in the production of, and does not endorse, this material. \textcopyright{} Integration Academy.}}
\fancyfoot[R]{\scriptsize\thepage}
\renewcommand{\headrulewidth}{1.2pt}
\renewcommand{\headrule}{\hbox to\headwidth{\color{gold}\leaders\hrule height \headrulewidth\hfill}}
\renewcommand{\footrulewidth}{0pt}
\setlength{\parindent}{0pt}
\setlength{\parskip}{4pt}

% Title block: \wstitle{code}{title}{unit line}
\newcommand{\wstitle}[3]{\begin{center}{\small\bfseries\color{gold}WORKSHEET #1}\\[3pt]{\LARGE\bfseries\color{navy}#2}\\[3pt]{\small\color{black!60}#3}\end{center}\vspace{2pt}}
% Section heading
\newcommand{\sech}[1]{\par\vspace{10pt}{\large\bfseries\color{navy}#1}\par\nobreak\vspace{2pt}{\color{rule}\hrule height0.8pt}\vspace{7pt}}
% A level tag printed beside a group of questions
\newcommand{\level}[1]{\par\vspace{6pt}{\footnotesize\bfseries\color{gold}\MakeUppercase{#1}}\par\nobreak\vspace{2pt}}

\newtcolorbox{ideabox}{colframe=ink,colback=inkbg,coltitle=white,colbacktitle=ink,fonttitle=\bfseries\small,title={Key ideas},arc=2pt,boxrule=0.8pt,left=8pt,right=8pt}
\newtcolorbox{example}[1]{colframe=rule,colback=white,coltitle=navy,colbacktitle=inkbg,fonttitle=\bfseries\small,title={#1},arc=2pt,boxrule=0.8pt,borderline west={3pt}{0pt}{ink},left=10pt,right=8pt}
\newtcolorbox{frbox}[1]{colframe=gold,colback=goldbg,coltitle=white,colbacktitle=gold,fonttitle=\bfseries\small,title={#1},arc=2pt,boxrule=0.8pt,breakable,left=8pt,right=8pt}
\newtcolorbox{solbox}{colframe=okgreen,colback=okbg,coltitle=white,colbacktitle=okgreen,fonttitle=\bfseries\small,title={Worked solutions},arc=2pt,boxrule=0.8pt,breakable,left=8pt,right=8pt}
\newtcolorbox{scorebox}[1]{colframe=navy,colback=white,coltitle=white,colbacktitle=navy,fonttitle=\bfseries\small,title={#1},arc=2pt,boxrule=0.8pt,breakable,left=8pt,right=8pt}

\newcommand{\soln}{\par\smallskip\textbf{\color{ink}Solution.}\ }
\newcommand{\result}[1]{\par\smallskip\textbf{#1}}
% blank working space of a given height
\newcommand{\work}[1]{\par\vspace{#1}}
% a point value in the margin of a scoring guide
\newcommand{\pts}[1]{\hfill{\footnotesize\bfseries\color{gold}[#1]}}

% Coordinate grid for graph questions: \begin{apgraph}{xmin}{xmax}{ymin}{ymax} ... \end{apgraph}
\newenvironment{apgraph}[4]{\begin{center}\begin{tikzpicture}\begin{axis}[
  axis lines=middle, xmin=#1, xmax=#2, ymin=#3, ymax=#4,
  grid=both, grid style={black!12}, xtick distance=1, ytick distance=1,
  ticklabel style={font=\scriptsize}, xlabel={$x$}, ylabel={$y$},
  width=10cm, height=7.2cm, samples=80, clip=false,
  every axis plot/.append style={line width=1.1pt, color=ink}]}{\end{axis}\end{tikzpicture}\end{center}}
\newcommand{\closeddot}[2]{\addplot[only marks, mark=*, mark size=2.2pt, color=ink] coordinates {(#1,#2)};}
\newcommand{\opendot}[2]{\addplot[only marks, mark=*, mark size=2.2pt, color=ink, mark options={fill=white}] coordinates {(#1,#2)};}

% Multiple-choice options: four across, or stacked when the options are long.
\newcommand{\opts}[4]{\par\vspace{3pt}\begin{tabular}{@{}p{0.235\linewidth}p{0.235\linewidth}p{0.235\linewidth}p{0.235\linewidth}@{}}(A)\ #1 & (B)\ #2 & (C)\ #3 & (D)\ #4\end{tabular}\par\vspace{9pt}}
\newcommand{\optsv}[4]{\par\vspace{3pt}\begin{tabular}{@{}l@{\ \ }p{0.9\linewidth}@{}}(A) & #1\\(B) & #2\\(C) & #3\\(D) & #4\end{tabular}\par\vspace{9pt}}

\fancyhead[R]{\footnotesize\color{black!60}Worksheet 7.1}
\begin{document}

\wstitle{7.1}{Writing a Differential Equation from a Description}{Unit 7: Differential Equations}

\begin{ideabox}
\begin{itemize}[leftmargin=*,itemsep=2pt,topsep=2pt]
  \item \textquotedblleft{}The rate of change of $y$\textquotedblright{} is a derivative: $\dfrac{dy}{dt}$.
  \item \textquotedblleft{}Proportional to\textquotedblright{} means a constant times: $ky$. \textquotedblleft{}Inversely proportional to\textquotedblright{} means a constant divided by: $\dfrac ky$.
  \item When something is added at one rate and removed at another, the equation is (rate in) $-$ (rate out).
  \item A given rate at a given value fixes the constant $k$. A decreasing quantity has a negative rate.
\end{itemize}
\end{ideabox}

\sech{Worked examples}
\begin{example}{Example 1 \textperiodcentered{} A direct translation}
The rate of change of $y$ with respect to $t$ is proportional to the square of $y$. Write a differential equation.\soln Rate: $\dfrac{dy}{dt}$. Proportional to $y^2$: $ky^2$.\result{$\dfrac{dy}{dt}=ky^2$.}
\end{example}

\begin{example}{Example 2 \textperiodcentered{} In and out}
A savings account with balance $B$ dollars earns interest at 3\% per year, and 2000 dollars per year is withdrawn. Write a differential equation, and find the balance that stays constant.\soln Rate in $0.03B$; rate out 2000.\result{$\dfrac{dB}{dt}=0.03B-2000$; constant when $B\approx66\,667$ dollars.}
\end{example}

\begin{example}{Example 3 \textperiodcentered{} Finding the constant}
A frozen meal at temperature $T$ warms in a room at 21 $^\circ$C at a rate proportional to $21-T$. When $T=5$ it is warming at 1.6 $^\circ$C per minute. Write the differential equation.\soln $\dfrac{dT}{dt}=k(21-T)$, and $1.6=k(16)$.\result{$\dfrac{dT}{dt}=0.1(21-T)$.}
\end{example}

\newpage
\sech{Practice}

\level{Core}
\begin{enumerate}[leftmargin=*,itemsep=2pt]
  \item Write a differential equation: the rate of change of $y$ with respect to $x$ is proportional to $x$.\work{2.3cm}
  \item Write a differential equation: the rate of change of $P$ with respect to $t$ is inversely proportional to $P$.\work{2.3cm}
  \item Write a differential equation: a quantity $N$ decreases at a rate proportional to $N$. State the sign of your constant.\work{2.3cm}
  \item A quantity satisfies $\dfrac{dy}{dt}=ky$, and $\dfrac{dy}{dt}=12$ when $y=3$. Find $k$.\work{2.3cm}
\end{enumerate}

\level{Exam level}
\begin{enumerate}[leftmargin=*,itemsep=2pt,start=5]
  \item A pollutant enters a lake at 12 kg per day, and each day 5\% of the amount $Q$ present is removed. Write a differential equation for $Q$, and find the amount at which $Q$ is steady.\work{2.8cm}
  \item A population $P$ grows at a rate proportional to the product of $P$ and $2000-P$. Write a differential equation, and find the value of $P$ at which the growth is fastest.\work{2.8cm}
  \item The depth $h$ cm of water in a tank falls at a rate proportional to $\sqrt{h}$. When $h=16$ the depth is falling at 2 cm per minute. Write the differential equation, including the constant.\work{2.8cm}
  \item In $\dfrac{dM}{dt}=kM$, the mass $M$ is in grams and the time $t$ in hours. What are the units of $k$?\work{2.8cm}
\end{enumerate}

\begin{frbox}{Question 9 \textperiodcentered{} Free-response practice}
A trainee's score $S$ on a skills test, out of 100, increases with practice at a rate proportional to $100-S$, where $t$ is the number of weeks of practice. When $S=40$, the score is increasing at 3 points per week.
\begin{enumerate}[label=(\alph*),leftmargin=*,itemsep=2pt]
  \item Write a differential equation for $S$ in terms of a constant $k$.\work{2.6cm}
  \item Find the value of $k$.\work{2.6cm}
  \item Find the rate at which the score is increasing when $S=80$, and describe what the model predicts as $S$ approaches 100.\work{2.6cm}
\end{enumerate}
\end{frbox}

\newpage
\begin{solbox}
\begin{enumerate}[leftmargin=*,itemsep=4pt]
  \item $\dfrac{dy}{dx}=kx$.
  \item $\dfrac{dP}{dt}=\dfrac kP$.
  \item $\dfrac{dN}{dt}=-kN$ with $k>0$ (or $\dfrac{dN}{dt}=kN$ with $k<0$).
  \item $12=3k$, so $k=4$.
  \item $\dfrac{dQ}{dt}=12-0.05Q$, steady when $Q=240$ kg.
  \item $\dfrac{dP}{dt}=kP(2000-P)$; $P(2000-P)$ is greatest at $P=1000$.
  \item $\dfrac{dh}{dt}=-k\sqrt h$ with $-2=-k(4)$: $\dfrac{dh}{dt}=-0.5\sqrt h$.
  \item Grams per hour divided by grams: per hour.
  \item \begin{enumerate}[label=(\alph*),leftmargin=*,itemsep=3pt]
    \item $\dfrac{dS}{dt}=k(100-S)$.
    \item $3=k(60)$, so $k=0.05$.
    \item $0.05(20)=1$ point per week. As $S\to100$ the rate tends to 0: improvement slows and the score levels off below 100.
  \end{enumerate}
\end{enumerate}
\end{solbox}

\end{document}
